% Part: applied-modal-logics
% Chapter: epistemic-logic
% Section: relational-models

\documentclass[../../../include/open-logic-section]{subfiles}

\begin{document}

\olfileid{aml}{el}{rel}

\olsection{संबंधी प्रतिमाने}

ज्ञानविषयक तर्कशास्त्रांची मूलभूत चिन्हार्थमीमांसक
संकल्पना सामान्य मोडल तर्कशास्त्रांतील संकल्पनेसारखीच
आहे. संबंधी प्रतिमानांत अजूनही जगांचा संच आणि कोणती
!!{propositional
variable}s कोणत्या जगांत ``सत्य''
मानायची हे ठरवणारे मूल्यांकन असते. आपण फक्त एकाच
कर्त्याचा विचार करत असू, तर नेहमीप्रमाणे एकच
प्राप्यता संबंध असतो. मात्र अनेक कर्त्यांच्या ज्ञानविषयक
तर्कशास्त्रात एकाच प्राप्यता संबंधाऐवजी प्राप्यता
संबंधांचा संच असतो: कर्तृचिन्हांच्या~$G$ संचातील
प्रत्येक~$a$ साठी एक संबंध.

\emph{संबंधी प्रतिमानात} जगांचा संच असतो; प्रत्येक
कर्त्यासाठी असलेल्या द्विस्थानी प्राप्यता संबंधांनी
ही जगे संबंधित असतात. त्यात कोणती !!{propositional variable}s कोणत्या जगांत सत्य आहेत हे ठरवणारे
मूल्यांकनही असते.

\begin{defn}
  अनेक कर्त्यांच्या ज्ञानविषयक भाषेचे \emph{प्रतिमान}
  म्हणजे $\mModel{M} = \tuple{W, R, V}$ असे त्रिकूट;
  येथे
  \begin{enumerate}
  \item $W$ हा ``जगांचा'' रिकामा नसलेला संच आहे,
  \item प्रत्येक $a \in G$ साठी ${R}_a$ हा~$W$ वरील
  द्विस्थानी प्राप्यता संबंध आहे, आणि
  \item $V$ हे प्रत्येक !!{propositional
    variable}~$p$ ला संभाव्य जगांचा संच~$V(p)$ नेमून
    देणारे फलन आहे.
  \end{enumerate}
  $R_a ww'$ असेल तेव्हा $w'$ हे कर्ता a साठी~$w$
  पासून \emph{प्राप्य} आहे असे म्हणतो. $w \in V(p)$
  असेल तेव्हा $p$ हे~$w$ येथे \emph{सत्य} आहे असे म्हणतो.
\end{defn}

याची कार्यपद्धती सामान्य मोडल तर्कशास्त्रासारखीच
आहे; फक्त आणखी प्राप्यता संबंध जोडलेले असतात.
दिलेल्या कर्त्याचा प्राप्यता संबंध सामान्यतः त्या
कर्त्याच्या माहितीच्या स्थितीबद्दल काहीतरी दर्शवतो,
असा अर्थ आपण लावू. उदाहरणार्थ, $R_a ww'$ हे
अनेकदा $w'$ हे जग~$w$ येथे~$a$ कडे असलेल्या
माहितीशी सुसंगत आहे, असे व्यक्त करते असे मानतो.
दुसऱ्या शब्दांत, $w$ येथे त्या कर्त्याला $w$ आणि
जग~$w'$ यांमधील फरक ओळखता येत नाही.

\end{document}
